% ----------------------------------------------------------------------
%  2.1.5. Получение высокой концентрации точечных дефектов
% ----------------------------------------------------------------------
\subsection{Получение высокой концентрации точечных дефектов}\label{ssec:production}

\begin{orientir}
При комнатной температуре равновесных дефектов ничтожно мало. Чтобы
изучать дефекты (или чтобы понимать, что происходит с материалом в
реакторе), их нужно получить много. Для этого есть три способа:
\textbf{закалка}, \textbf{пластическая деформация} и \textbf{облучение}
частицами высокой энергии. У каждого свои возможности и ограничения~— их
мы и сравним.
\end{orientir}

\subsubsection{Закалка}

\paragraph{Идея}
При высокой температуре вакансии очень подвижны и быстро приходят в
равновесие со стоками: их концентрация равна равновесной
$C_v(T_{\text{з}})$. Если кристалл охладить так быстро, что вакансии не
успеют уйти на стоки, эта концентрация «заморозится». При низкой
температуре кристалл окажется \emph{пересыщенным} вакансиями.

Процедура: образец выдерживают при \term{температуре закалки}
$T_{\text{з}}$ до установления равновесия, затем быстро охлаждают до
низкой температуры $T_{\text{к}}$. Параметры выбирают так:
\begin{itemize}
\item $T_{\text{з}}$~— достаточно высокая, чтобы вакансий было много
(но ниже точки плавления);
\item $T_{\text{к}}$~— достаточно низкая, чтобы за время измерений
вакансии не успели уйти на стоки;
\item скорость охлаждения $V_{\text{з}}$~— достаточно большая, чтобы
вакансии не ушли на стоки во время охлаждения.
\end{itemize}
Заметный избыток вакансий даёт уже обычная закалка в воду.

\paragraph{Потери вакансий при охлаждении}
Пусть до стока вакансия должна в среднем сделать $n$ перескоков, а
частота её перескоков $\Gamma=\nu_0\exp(-\Evm/kT)$. Тогда избыток
вакансий убывает со скоростью
\[
\frac{\d C_v}{\d t}=-\frac{\Gamma}{n}\bigl(C_v-C_v^0(T)\bigr),
\qquad C_v^0(T)=\Ex{-\frac{\EvF}{kT}} .
\]
Пусть охлаждение экспоненциальное: $\d T/\d t=-\beta(T-T_{\text{к}})$,
где $\beta$~— константа охлаждения. Перейдя от времени к температуре,
получим уравнение
\[
\frac{\d C_v}{\d T}=\frac{\nu_0}{n\beta\,(T-T_{\text{к}})}\Ex{-\frac{\Evm}{kT}}
\bigl(C_v-C_v^0(T)\bigr),
\]
которое интегрируется численно (рис.~\ref{fig:quench}). Результат
определяется одним параметром $\gamma=\beta n/\nu_0$.

\begin{figure}[ht]
\centering
\begin{tikzpicture}
\begin{axis}[width=0.7\textwidth, height=6.2cm, ymode=log, x dir=reverse,
  xmin=300, xmax=1300, ymin=1e-8, ymax=2,
  xlabel={$T$, К (охлаждение $\rightarrow$)}, ylabel={$C_v/C_v^0(T_{\text{з}})$},
  ytick={1e-8,1e-6,1e-4,1e-2,1}, grid=major, grid style={fmGray!20},
  legend pos=south west, legend cell align=left]
  \addplot[black, dashed, thick, domain=520:1300, samples=60] {exp(-1.1/(8.617e-5*x)+1.1/(8.617e-5*1300))};
  \addlegendentry{равновесие $C_v^0(T)$}
  \addplot[fmBlue, very thick] coordinates {(1300,1) (1250,0.9986) (1150,0.9918) (1050,0.9855) (950,0.9815) (850,0.9795) (750,0.9788) (650,0.9785) (300,0.9785)};
  \addlegendentry{$\gamma=10^{-2}$}
  \addplot[fmGreen, very thick] coordinates {(1300,1) (1250,0.8898) (1200,0.7125) (1150,0.5552) (1100,0.4341) (1050,0.3466) (1000,0.2861) (950,0.2455) (900,0.2193) (850,0.2028) (800,0.193) (750,0.1873) (700,0.1843) (600,0.1822) (300,0.1818)};
  \addlegendentry{$\gamma=10^{-4}$}
  \addplot[fmRed, very thick] coordinates {(1300,1) (1250,0.6789) (1200,0.4444) (1150,0.2806) (1100,0.1701) (1050,0.09851) (1000,0.05419) (950,0.02816) (900,0.01376) (850,0.006322) (800,0.002761) (750,0.001192) (700,0.0005548) (650,0.0003141) (600,0.0002283) (550,0.0001995) (500,0.0001908) (400,0.0001881) (300,0.0001881)};
  \addlegendentry{$\gamma=10^{-6}$}
\end{axis}
\end{tikzpicture}
\caption{Концентрация вакансий при закалке с 1300~К (расчёт для
$\EvF=1{,}1\eV$, $\Evm=0{,}7\eV$). Пунктир~— равновесная концентрация.
При большом $\gamma$ почти все вакансии «замораживаются»; при малом
кристалл успевает следить за равновесием до низких температур, и избыток
теряется}
\label{fig:quench}
\end{figure}

\begin{glavnoe}
Закалка тем эффективнее, чем больше $\gamma=\beta n/\nu_0$, т.\,е. чем
быстрее охлаждение (больше $\beta$) и чем меньше стоков (больше $n$~—
число перескоков до стока).
\end{glavnoe}

\paragraph{Ограничения закалки}
\begin{enumerate}[label=\arabic*)]
\item \emph{Термические напряжения.} Поверхность и сердцевина образца
остывают с разной скоростью; возникающие напряжения могут вызвать
движение дислокаций, а движущиеся дислокации сами рождают точечные
дефекты, уже не «термодинамическим» путём (см. ниже). Это искажает
результат. Ван~Бюрен получил предельную скорость закалки проволоки
радиусом $r$, при которой пластической деформации ещё нет:
\begin{equation}\label{eq:vanburen}
V<V_{\max}=\frac{5(1-\nu)\,\delta\,\tau_{\text{кр}}}{\alpha_TGr^2},
\qquad
r<r_{\max}=\sqrt{\frac{5(1-\nu)\,\delta\,\tau_{\text{кр}}}{\alpha_TGV}},
\end{equation}
где $\nu$~— коэффициент Пуассона, $\delta$~— температуропроводность,
$\alpha_T$~— коэффициент теплового расширения, $\tau_{\text{кр}}$~—
критическое напряжение сдвига, $G$~— модуль сдвига
(табл.~\ref{tab:vanburen}).
\item \emph{Образование комплексов.} Во время закалки кристалл пересыщен,
и вакансии объединяются. Энергия образования дивакансии
$E_{2v}^{\mathrm F}=2\EvF-B_2$ больше, чем у вакансии, поэтому
концентрация дивакансий растёт с температурой быстрее и при высоких
$T_{\text{з}}$ приближается к концентрации моновакансий. Расчёт показывает:
если $B_2>0{,}3\eV$, более половины вакансий за время закалки
превращаются в дивакансии. При закалке с очень высоких температур
вакансии успевают собраться в диски диаметром $\sim10$~нм, которые
«захлопываются» в дислокационные петли (п.~\ref{ssec:disl-origin}).
\item \emph{Только вакансии.} Закалка не даёт межузельных атомов, а
концентрация вакансий ограничена значением у точки плавления
($\sim10^{-4}$).
\end{enumerate}

\begin{table}[ht]
\centering
\caption{Предельные скорости закалки и радиусы проволоки, при которых
ещё нет пластической деформации}
\label{tab:vanburen}
\small
\begin{tabular}{@{}lccccccc@{}}
\toprule
 & Cu & Au & Ag & Al & Ni & Fe & W \\
\midrule
$V_{\max}$, $10^4$~К/с, при $r=0{,}1$~мм & 1,7 & 2,4 & 2,4 & 1,0 & 0,8 & 0,6 & 0,11 \\
$r_{\max}$, мм, при $V=10^4$~К/с & 0,13 & 0,16 & 0,16 & 0,10 & 0,09 & 0,08 & 0,033 \\
\bottomrule
\end{tabular}
\end{table}
\noindent Из~\eqref{eq:vanburen} видно, что $r_{\max}=0{,}1\text{~мм}\cdot
\sqrt{V_{\max}(0{,}1\text{~мм})/10^4}$; именно так связаны строки
таблицы\ispr{В табл.~2.3 оригинала для W указано $V_{\max}=1{,}1$; это не
согласуется с $r_{\max}=0{,}033$~мм в той же строке. Согласованное значение
$V_{\max}\approx0{,}11\cdot10^4$~К/с. В формуле для $r_{\max}$ в оригинале
потерян квадратный корень.}.

\subsubsection{Пластическая деформация}

При деформации дислокации движутся и размножаются (п.~\ref{sec:disl}), а
движущиеся дислокации порождают точечные дефекты. Простейший пример~—
встреча двух краевых дислокаций противоположного знака
(рис.~\ref{fig:annih}). Краевая дислокация~— это край лишней атомной
полуплоскости; у дислокации одного знака полуплоскость вставлена сверху,
у другой~— снизу.

\begin{itemize}
\item Если дислокации скользят по \emph{одной} плоскости, то при встрече
две полуплоскости сливаются в одну целую плоскость~— дислокации
аннигилируют, и остаётся совершенный кристалл (рис.~\ref{fig:annih},\,а).
\item Если их плоскости скольжения разнесены так, что между краями
полуплоскостей остаётся пустой промежуток в один атом, получается целая
плоскость, в которой не хватает ряда атомов,~— \emph{цепочка вакансий}
(рис.~\ref{fig:annih},\,б).
\item Если края полуплоскостей «перекрываются», в получившейся плоскости
оказывается лишний ряд атомов~— \emph{цепочка межузельных атомов}
(рис.~\ref{fig:annih},\,в).
\end{itemize}

\begin{figure}[ht]
\centering
\begin{tikzpicture}[scale=0.5]
\newcommand{\annpanel}[4]{% #1 сдвиг, #2 верх. край, #3 ниж. край, #4 метка
  \begin{scope}[shift={(#1,0)}]
  \foreach \x in {-2,-1,1,2}{\foreach \y in {-3,...,3}{\node[atom=3.4mm] at (\x,\y) {};}}
  \foreach \y in {#2,...,3}{\node[atomB=3.4mm] at (0,\y) {};}
  \foreach \y in {-3,...,#3}{\node[atomB=3.4mm, fill=atomB!60] at (0,\y) {};}
  \node[lbl] at (0,-4.3) {#4};
  \end{scope}}
\annpanel{0}{0}{-1}{а}
\annpanel{7}{1}{-1}{б}
\node[vac=3.4mm] at (7,0) {};
\annpanel{14}{1}{-1}{в}
\node[atomB=3.4mm] at (14.0,0.3) {};
\node[atomI=3.4mm] at (14.0,-0.3) {};
\end{tikzpicture}
\caption{Результат встречи краевых дислокаций противоположного знака.
Тёмно-оранжевые атомы~— бывшая полуплоскость, вставленная сверху,
светло-оранжевые~— снизу. \textit{а}~— полуплоскости сливаются точно:
совершенная решётка; \textit{б}~— между ними остаётся промежуток: ряд
вакансий (перпендикулярно рисунку); \textit{в}~— полуплоскости
перекрываются: на одно место претендуют два атома~— ряд межузельных
атомов (гантелей)}
\label{fig:annih}
\end{figure}

Важно, что при деформации энергию на образование дефектов дают
\emph{внешние силы}, поэтому, в отличие от теплового равновесия,
межузельные атомы рождаются наравне с вакансиями. Точечные дефекты
образуются также при движении дислокаций с порогами
(п.~\ref{ssec:disl-motion}). Полуколичественные оценки дают долю
дефектов, пропорциональную степени деформации, с коэффициентом
$10^{-4}$--$10^{-5}$. Цепочки вакансий могут затем распадаться
(диспергировать) на отдельные дефекты; вероятно, в итоге преобладают
дивакансии и более крупные комплексы. Дислокации служат и стоками для
точечных дефектов.

\paragraph{Ограничения}
Деформация даёт большие концентрации и вакансий, и межузельных атомов,
но одновременно меняется число и расположение дислокаций. Дислокации
взаимодействуют с точечными дефектами, поэтому «окружение» дефектов
сильно отличается от отожжённого кристалла, и опыты трудно истолковать.
К тому же эффективность и закалки, и деформации сильно зависит от
температуры.

\subsubsection{Облучение частицами высокой энергии}

Облучение свободно от многих недостатков первых двух методов: оно даёт
высокие концентрации и вакансий, и межузельных атомов практически
независимо от температуры. Для студентов ядерных специальностей этот
способ особенно важен: именно так возникают дефекты в материалах
реакторов.

\paragraph{Пороговая энергия смещения}
Заряженные частицы при прохождении через кристалл в основном ионизуют и
возбуждают электроны. Дефекты же создаются \emph{упругими} соударениями с
ядрами атомов. Если атом решётки получил энергию $T$, большую
\term{пороговой энергии смещения} $E_d$, он покидает узел, проходит
несколько межатомных расстояний и останавливается в междоузлии.
Образуется пара Френкеля.

\begin{glavnoe}
$E_d\approx10$--$30\eV$ (типично $\approx25\eV$)~— в несколько раз
больше, чем $3$--$6\eV$, нужные для создания пары Френкеля обратимым,
термодинамическим путём.
\end{glavnoe}

\begin{smysl}[Почему порог так высок (оценка Зейтца)]
Энергия связи $E_c$ у самых прочных кристаллов $5$--$6\eV$. При
сублимации атом уходит с поверхности, где у него половина связей;
значит, чтобы \emph{обратимо} вынуть атом из объёма, нужно $\approx2E_c$.
При ударе атом выбивается резко и \emph{необратимо}: соседи не успевают
подстроиться, и требуется по меньшей мере $\approx4E_c\approx20$--$30\eV$.
На самом деле $E_d$ сильно зависит от направления удара относительно
кристаллографических осей.
\end{smysl}

\paragraph{Сколько энергии может передать частица}
Энергия, переданная атому, меняется от нуля (скользящий удар) до
максимальной $T_{\max}$ (лобовой удар). Из законов сохранения энергии и
импульса для нерелятивистской частицы массы $m$ с энергией $E$ и атома
массы $M$:
\begin{keyeq}\label{eq:Tmax}
T_{\max}=\frac{4mM}{(m+M)^2}\,E .
\end{keyeq}
Чем сильнее различаются массы, тем меньшую долю энергии частица может
передать. Для электронов (они релятивистские уже при $\sim1$~МэВ)
$T_{\max}=2E(E+2m_ec^2)/(Mc^2)$. Приравнивая $T_{\max}=E_d$, находим
пороговую энергию частицы (табл.~\ref{tab:thresholds}).

\begin{table}[ht]
\centering
\caption{Минимальная энергия бомбардирующей частицы, при которой
возникают смещения ($E_d=25\eV$)}
\label{tab:thresholds}
\small
\begin{tabular}{@{}lcccc@{}}
\toprule
Частица & \multicolumn{4}{c}{Атомная масса мишени} \\
\cmidrule(l){2-5}
 & 10 & 50 & 100 & 200 \\
\midrule
Нейтроны, протоны, эВ & 70 & 325 & 638 & 1263 \\
$\alpha$-частицы, эВ & 31 & 91 & 169 & 325 \\
Электроны, $\gamma$-кванты, МэВ & 0,10 & 0,41 & 0,68 & 1,10 \\
\bottomrule
\end{tabular}
\end{table}
\noindent Строки для нейтронов и $\alpha$-частиц легко проверить
по~\eqref{eq:Tmax}: например, для нейтрона и мишени с $M=50$ доля
$4\cdot50/51^2=0{,}077$, и $E=25/0{,}077=325\eV$\ispr{В табл.~2.4
оригинала для $\alpha$-частиц и мишени массой 200 указано «3,25»;
расчёт по~\eqref{eq:Tmax} даёт 325~эВ.}.

\paragraph{Каскады смещений}
Если переданная энергия ненамного больше $E_d$, образуется одна пара
Френкеля. Если же \term{первично выбитый атом} (ПВА) получил энергию,
намного большую $E_d$, он сам выбивает другие атомы, те~— следующие, и
развивается \term{каскад смещений} (рис.~\ref{fig:cascade}). Его
структура неоднородна:
\begin{itemize}
\item межузельные атомы уносятся от центра каскада на заметное расстояние
\emph{цепочками замещающих соударений} вдоль плотноупакованных рядов
(атом выбивает соседа, занимает его место, тот~— следующего, и т.\,д.);
\item у вакансий такого динамического механизма переноса нет, и они
остаются в центре каскада компактной \emph{обеднённой зоной}.
\end{itemize}
В итоге получаются изолированные межузельные атомы (в гантельной
конфигурации) и различные комплексы вакансий.

\begin{figure}[ht]
\centering
\begin{tikzpicture}[scale=0.42]
  \foreach \x in {-9,...,9}{\foreach \y in {-5,...,5}{
     \pgfmathsetmacro{\rr}{(\x/1.3)^2+(\y)^2}
     \pgfmathparse{(\rr<5.5 && mod(abs(\x+3*\y),3)!=0) ? 1 : 0}
     \ifnum\pgfmathresult=1 \node[vac=3mm] at (\x,\y) {};
     \else \node[atom=3mm, fill=atomA!70] at (\x,\y) {};\fi}}
  \foreach \p in {(8.5,0.5),(-8.5,-1.5),(3.5,4.5),(-4.5,-4.5),(6.5,-3.5),(-6.5,3.5),(0.5,-5.5)}{
     \node[atomI=3mm] at \p {};}
  \draw[arrB] (-12,3) -- (-6.3,1.2);
  \node[lbl, fmRed] at (-12.5,3.8) {ПВА};
  \foreach \ax/\ay/\bx/\by in {1.5/0/8/0.5, -2.5/-0.6/-8/-1.4, 0.6/1.6/3.3/4.1,
      -1.4/-2/-4.2/-4.1, 2.4/-1.4/6.2/-3.3, -2.6/1.4/-6.2/3.3, 0.2/-2.4/0.45/-5.1}{
     \draw[-{Stealth[length=1.6mm]}, fmRed!80, line width=0.5pt, dashed] (\ax,\ay) -- (\bx,\by);}
  \draw[fmBlue, thick, dashed] (0,0) ellipse (3.2 and 2.5);
  \node[lbl, fmBlue, align=center] at (0,-7) {обеднённая (вакансионная) зона};
\end{tikzpicture}
\caption{Схема каскада смещений: в центре~— зона, обеднённая атомами
(пустые кружки~— вакансии); межузельные атомы (красные) унесены на
периферию цепочками замещающих соударений (пунктирные стрелки)}
\label{fig:cascade}
\end{figure}

\paragraph{Число смещённых атомов}
Среднее число смещённых атомов в единице объёма
\begin{equation}\label{eq:Nd}
N_d=\varphi\,t\,n_0\,\sigma_d\,\nu ,
\end{equation}
где $\varphi$~— плотность потока частиц, $t$~— время облучения, $n_0$~—
число атомов в единице объёма, $\sigma_d$~— сечение столкновений,
приводящих к смещению (в расчёте на атом), $\nu$~— среднее число смещений
на один ПВА. Последнее находят по общепринятой модели
Норджетта--Робинсона--Торренса (стандарт NRT\ispr{В оригинале модель
названа «Торренса--Робинсона--Норгетта (TRN)»; общепринятое название~—
стандарт NRT (Norgett, Robinson, Torrens, 1975).}):
\begin{equation}\label{eq:NRT}
\nu=\frac{\chi\,\bar T}{2E_d},\qquad\chi\approx0{,}8,
\end{equation}
где $\bar T$~— средняя энергия, переданная ПВА и затраченная на
упругие соударения.

\begin{table}[ht]
\centering
\caption{Среднее число смещённых атомов на одно первичное соударение
при облучении нейтронами с энергией 1~МэВ}
\label{tab:nu}
\small
\begin{tabular}{@{}lcccc@{}}
\toprule
 & Железо & Медь & Германий\footnotemark & Золото \\
\midrule
Атомная масса & 56 & 64 & 73 & 197 \\
$\nu$ & 390 & 380 & 290 & 140 \\
\bottomrule
\end{tabular}
\end{table}
\footnotetext{\textcolor{fmRed}{\sffamily\bfseries Исправление.}~В табл.~2.5 оригинала атомная
масса германия указана как 93; правильно~— 72,6.}

\begin{primer}[Смещения на атом в реакторе]
Отношение $N_d/n_0=\varphi t\sigma_d\nu$~— это число смещений,
пришедшихся в среднем на каждый атом, \term{сна} (\emph{смещений на
атом}, англ.\ \emph{dpa}). Это основная мера радиационного повреждения
реакторных материалов. Для железа в активной зоне быстрого реактора:
$\varphi\approx10^{15}$~нейтр./(см$^2\cdot$с), $\sigma_d\approx3\cdot10^{-24}$~см$^2$,
$\nu\approx390$:
\[
\frac{\d(\text{сна})}{\d t}\approx10^{15}\cdot3\cdot10^{-24}\cdot390\approx
1{,}2\cdot10^{-6}~\text{с}^{-1}\approx 35~\text{сна в год}.
\]
За год каждый атом покидает своё место в среднем несколько десятков раз!
Большая часть пар Френкеля тут же рекомбинирует, но выжившие дефекты
определяют распухание, упрочнение и охрупчивание материалов под
облучением.
\end{primer}

\paragraph{Электроны и $\gamma$-кванты}
При облучении электронами с энергией 1--3~МэВ $\nu\approx1$: на каждое
первичное соударение приходится одна, самое большее несколько пар
Френкеля. ПВА не хватает энергии на каскад, поэтому электронное облучение
даёт \emph{самое простое} распределение дефектов~— изолированные пары.
Так же действует $\gamma$-облучение (смещения создают комптоновские
электроны и фотоэлектроны). Оно гораздо менее эффективно, чем
бомбардировка электронами, зато $\gamma$-кванты проникают глубоко и
повреждают образец однородно.

\begin{table}[ht]
\centering
\caption{Сравнение способов получения точечных дефектов}
\label{tab:methods-compare}
\small
\begin{tabularx}{\textwidth}{@{}l>{\raggedright\arraybackslash}X>{\raggedright\arraybackslash}X@{}}
\toprule
Способ & Что получаем & Ограничения \\
\midrule
Закалка & только вакансии и их комплексы; $C_v\lesssim10^{-4}$ &
  потери на стоках; термические напряжения; образование комплексов \\
Деформация & вакансии и межузельные атомы, большие концентрации &
  одновременно меняется дислокационная структура; трудно истолковать \\
Облучение & пары Френкеля, каскады; высокие концентрации при любой $T$ &
  сложное пространственное распределение (каскады); нужна установка \\
\bottomrule
\end{tabularx}
\end{table}

\begin{itogi}
\begin{itemize}
\item Закалка замораживает высокотемпературную концентрацию вакансий;
эффективность растёт с $\gamma=\beta n/\nu_0$. Межузельных атомов не
даёт; мешают термические напряжения и образование комплексов.
\item Пластическая деформация рождает вакансии и межузельные атомы
(например, при аннигиляции краевых дислокаций на соседних плоскостях),
пропорционально степени деформации.
\item Облучение: смещение при $T>E_d\approx25\eV$;
$T_{\max}=4mME/(m+M)^2$; тяжёлые частицы создают каскады, электроны~—
изолированные пары Френкеля. Мера повреждения~— сна (dpa).
\end{itemize}
\end{itogi}

\begin{proverka}
\item Почему закалкой нельзя получить межузельные атомы?
\item Как повлияет на эффективность закалки уменьшение размера зерна?
\item Какую максимальную энергию передаст нейтрон с энергией 1~МэВ ядру
железа ($M=56$)?
\item Почему вакансии в каскаде остаются в центре, а межузельные атомы
уходят на периферию?
\end{proverka}
